\chapter{§3.9 坐标几何}
\label{sec:3.9}


\prerequisite{§2.10, §3.1}

\gradelevel{8 年级}


\section*{平面直角坐标系}


\begin{explore}


电影院里，你的座位票上写着" $5$ 排 $12$ 号"。这两个数字精确地确定了你在整个影厅中的位置。数学中也有类似的方法——用两个数来确定平面上一个点的位置。

\end{explore}

\begin{understand}


在平面上画两条互相垂直且有公共原点的数轴，就构成了\textbf{平面直角坐标系}（笛卡尔坐标系）。

\begin{figure}[htbp]
\centering
\begin{tikzpicture}[every node/.style={font=\small}, scale=0.8]
  % Axes
  \draw[ray arrow] (-3.3,0) -- (3.5,0) node[right] {$x$};
  \draw[ray arrow] (0,-2.8) -- (0,3.3) node[above] {$y$};
  \node[below left] at (0,0) {$O$};
  % Grid ticks
  \foreach \x in {-3,-2,-1,1,2,3} {
    \draw (\x, 0.07) -- (\x,-0.07) node[below, font=\scriptsize] {\x};
  }
  \foreach \y in {-2,-1,1,2,3} {
    \draw (0.07,\y) -- (-0.07,\y) node[left, font=\scriptsize] {\y};
  }
  % Sample points
  \fill[draw=red!70!black, fill=red!70!black] (2,3) circle (2.5pt)
    node[above right] {$P(2,3)$};
  \fill[draw=blue!70!black, fill=blue!70!black] (-2,1) circle (2.5pt)
    node[above left] {$Q(-2,1)$};
  \fill[draw=green!60!black, fill=green!60!black] (1,-2) circle (2.5pt)
    node[below right] {$R(1,-2)$};
  % Dashed projection for P
  \draw[aux line] (2,0) -- (2,3) -- (0,3);
  % Quadrant labels
  \node[font=\scriptsize, gray] at (2,1.5) {第一象限};
  \node[font=\scriptsize, gray] at (-2,1.5) {第二象限};
  \node[font=\scriptsize, gray] at (-2,-1.5) {第三象限};
  \node[font=\scriptsize, gray] at (2,-1.5) {第四象限};
\end{tikzpicture}
\caption{直角坐标系：点 $P(x,y)$ 由横坐标 $x$ 和纵坐标 $y$ 唯一确定}
\label{fig:cartesian-plane}
\end{figure}


\begin{itemize}
  \item \textbf{水平的数轴}叫做 $x$ 轴（横轴），向右为正方向。
  \item \textbf{竖直的数轴}叫做 $y$ 轴（纵轴），向上为正方向。
  \item 两轴的交点 $O$ 叫做\textbf{原点}。
\end{itemize}


坐标系把平面分成四个\textbf{象限}：

\begin{center}
\begin{tabular}{lll}
\toprule
象限 & $x$ 的符号 & $y$ 的符号 \\
\midrule
第一象限 & $+$ & $+$ \\
第二象限 & $-$ & $+$ \\
第三象限 & $-$ & $-$ \\
第四象限 & $+$ & $-$ \\
\bottomrule
\end{tabular}
\end{center}


\end{understand}

\begin{quote}
坐标轴上的点不属于任何象限。
\end{quote}


\begin{workedexamples}


\textbf{例 1.} 判断下列点在哪个象限或坐标轴上： $A(3, 2)$ 、 $B(-1, 4)$ 、 $C(0, -5)$ 、 $D(-3, -2)$ 。

\textbf{解：}
$A(3, 2)$ ： $x > 0$ ， $y > 0$ ，第一象限。

$B(-1, 4)$ ： $x < 0$ ， $y > 0$ ，第二象限。

$C(0, -5)$ ： $x = 0$ ，在 $y$ 轴上（不属于任何象限）。

$D(-3, -2)$ ： $x < 0$ ， $y < 0$ ，第三象限。

\textbf{例 2.} 若点 $P(a, b)$ 在第四象限，则 $a$ 和 $b$ 分别满足什么条件？

\textbf{解：}
第四象限： $a > 0$ ， $b < 0$ 。

\textbf{例 3.} 若点 $P(m-1, m+2)$ 在 $y$ 轴上，求 $m$ 的值和点 $P$ 的坐标。

\textbf{解：}
在 $y$ 轴上的点横坐标为 $0$ ： $m - 1 = 0$ ， $m = 1$ 。

$P = (0, 3)$ 。

\end{workedexamples}

\begin{keytakeaway}


\begin{itemize}
  \item 平面直角坐标系用有序数对 $(x, y)$ 表示点的位置。
  \item 四个象限由两条坐标轴划分，坐标轴上的点不属于任何象限。
  \item $x$ 轴上的点纵坐标为 $0$ ， $y$ 轴上的点横坐标为 $0$ 。
\end{itemize}


\end{keytakeaway}

\begin{practice}


\begin{enumerate}
  \item 写出下列点所在的象限或坐标轴： $A(5, -3)$ 、 $B(-4, 0)$ 、 $C(-2, -7)$ 、 $D(0, 0)$ 。
  \item 若点 $P(2a-4, a+1)$ 在第二象限，求 $a$ 的取值范围。
\end{enumerate}


\end{practice}



\section*{点的坐标}


\begin{explore}


在方格纸上画一个坐标系，标出几个格点（格子的交叉点），读出它们的坐标——这就是"点的坐标"最基本的操作。反过来，给定坐标 $(3, -2)$ ，你能准确地在方格纸上找到这个点。

\end{explore}

\begin{understand}


\textbf{坐标的读取}：从点向 $x$ 轴作垂线，垂足在 $x$ 轴上的读数即为\textbf{横坐标}（ $x$ ）；从点向 $y$ 轴作垂线，垂足在 $y$ 轴上的读数即为\textbf{纵坐标}（ $y$ ）。

特殊点的坐标：
\begin{itemize}
  \item 原点： $(0, 0)$
  \item $x$ 轴上的点： $(a, 0)$
  \item $y$ 轴上的点： $(0, b)$
\end{itemize}


\textbf{坐标与几何位置的关系}：
\begin{itemize}
  \item 横坐标相同的两个点在同一条竖直线上。
  \item 纵坐标相同的两个点在同一条水平线上。
\end{itemize}


\end{understand}

\begin{workedexamples}


\textbf{例 1.} 在坐标平面上画出 $A(2, 3)$ 、 $B(-1, 2)$ 、 $C(3, -1)$ 、 $D(-2, -3)$ ，并说明它们分别在哪个象限。

\textbf{解：}
$A(2, 3)$ ：从原点向右 $2$ 、向上 $3$ ，第一象限。

$B(-1, 2)$ ：从原点向左 $1$ 、向上 $2$ ，第二象限。

$C(3, -1)$ ：从原点向右 $3$ 、向下 $1$ ，第四象限。

$D(-2, -3)$ ：从原点向左 $2$ 、向下 $3$ ，第三象限。

\textbf{例 2.} 已知 $A(-3, 0)$ 、 $B(0, 4)$ ，求线段 $AB$ 的长度。

\textbf{解：}
$AB = \sqrt{(-3-0)^2 + (0-4)^2} = \sqrt{9 + 16} = \sqrt{25} = 5$ 。

（距离公式将在后面正式学习，这里也可以看出 $\triangle AOB$ 是直角三角形，直角在原点，两直角边分别为 $3$ 和 $4$ ，斜边 $= 5$ 。）

\textbf{例 3.} 正方形 $ABCD$ 的顶点 $A(0,0)$ 、 $B(4,0)$ 。 $C$ 、 $D$ 在 $x$ 轴上方，求 $C$ 、 $D$ 的坐标。

\textbf{解：}
$AB = 4$ ，沿 $x$ 轴。正方形的边长 $= 4$ 。

$D$ 在 $A$ 正上方（或按逆时针方向）， $D = (0, 4)$ 。

$C$ 在 $B$ 正上方， $C = (4, 4)$ 。

\end{workedexamples}

\begin{keytakeaway}


\begin{itemize}
  \item 坐标 $(x, y)$ ： $x$ 是横坐标， $y$ 是纵坐标。
  \item 坐标的本质是用数来描述位置。
\end{itemize}


\end{keytakeaway}

\begin{practice}


\begin{enumerate}
  \item 在坐标平面上标出以下各点： $A(0, 4)$ 、 $B(3, 0)$ 、 $C(-2, 3)$ 、 $D(5, -2)$ 。
  \item 等腰直角三角形 $OAB$ 中， $O$ 为原点， $A(6, 0)$ 。 $B$ 在 $y$ 轴正半轴上，求 $B$ 的坐标。
\end{enumerate}


\end{practice}



\section*{用坐标表示平移}


\begin{explore}


电子游戏中，角色向右走 $3$ 步、向上跳 $2$ 步——角色的坐标就从 $(x, y)$ 变成了 $(x+3, y+2)$ 。平移在坐标中有非常简洁的表达。

\end{explore}

\begin{understand}


将图形向右平移 $a$ 个单位、向上平移 $b$ 个单位，则每个点的坐标变化为：

\[
(x, y) \to (x + a, y + b)
\]


向左平移 $a$ 个单位等同于 $a$ 取负值，向下平移 $b$ 个单位等同于 $b$ 取负值。

\end{understand}

\begin{workedexamples}


\textbf{例 1.} $\triangle ABC$ 的顶点为 $A(1, 2)$ 、 $B(4, 1)$ 、 $C(3, 5)$ 。将三角形向左平移 $3$ 个单位、向下平移 $2$ 个单位，求新的顶点坐标。

\textbf{解：}
向左 $3$ ，向下 $2$ ： $(x, y) \to (x - 3, y - 2)$ 。

$A' = (1 - 3, 2 - 2) = (-2, 0)$

$B' = (4 - 3, 1 - 2) = (1, -1)$

$C' = (3 - 3, 5 - 2) = (0, 3)$

\textbf{例 2.} 点 $P(2, 5)$ 经平移后到达 $P'(-1, 3)$ 。求平移的向量（方向和距离）。

\textbf{解：}
水平移动： $-1 - 2 = -3$ （向左 $3$ ）。

竖直移动： $3 - 5 = -2$ （向下 $2$ ）。

同一次平移将 $Q(4, -1)$ 变换为 $Q'(4-3, -1-2) = (1, -3)$ 。

\textbf{例 3.} 直线 $y = 2x + 1$ 向上平移 $3$ 个单位后的方程是什么？

\textbf{解：}
向上平移 $3$ 个单位：用 $y - 3$ 代替 $y$ （或等价地， $y \to y - 3$ ）。

新方程： $y - 3 = 2x + 1$ ，即 $y = 2x + 4$ 。

\end{workedexamples}

\begin{keytakeaway}


\begin{itemize}
  \item 平移在坐标中表现为坐标的加减： $(x, y) \to (x + a, y + b)$ 。
  \item 平移不改变图形形状和大小，也不改变直线的斜率。
\end{itemize}


\end{keytakeaway}

\begin{practice}


\begin{enumerate}
  \item 点 $A(-3, 4)$ 向右平移 $5$ 、向下平移 $3$ 后的坐标。
  \item 线段 $AB$ 中， $A(1, 2)$ 、 $B(5, 6)$ 。将 $AB$ 向右平移 $3$ 个单位后，新线段的两端点坐标是什么？新线段的长度与原来相比如何？
\end{enumerate}


\end{practice}



\section*{用坐标表示轴对称}


\begin{explore}


镜子放在 $y$ 轴的位置，你站在 $(3, 2)$ ，你在镜中的像在哪里？答案是 $(-3, 2)$ ——横坐标变号，纵坐标不变。

\end{explore}

\begin{understand}


\textbf{关于 $x$ 轴对称}： $(x, y) \to (x, -y)$ （纵坐标变号）

\textbf{关于 $y$ 轴对称}： $(x, y) \to (-x, y)$ （横坐标变号）

\textbf{关于原点对称}： $(x, y) \to (-x, -y)$ （两个坐标都变号）

\begin{center}
\begin{tabular}{lll}
\toprule
对称类型 & 坐标变换 & 什么变什么不变 \\
\midrule
关于 $x$ 轴 & $(x, y) \to (x, -y)$ & $x$ 不变， $y$ 变号 \\
关于 $y$ 轴 & $(x, y) \to (-x, y)$ & $y$ 不变， $x$ 变号 \\
关于原点 & $(x, y) \to (-x, -y)$ & $x$ 、 $y$ 都变号 \\
\bottomrule
\end{tabular}
\end{center}


\end{understand}

\begin{workedexamples}


\textbf{例 1.} 点 $A(3, -4)$ 关于 $x$ 轴、 $y$ 轴、原点的对称点分别是什么？

\textbf{解：}
关于 $x$ 轴： $A_1 = (3, 4)$ 。

关于 $y$ 轴： $A_2 = (-3, -4)$ 。

关于原点： $A_3 = (-3, 4)$ 。

\textbf{例 2.} $\triangle ABC$ 的顶点为 $A(1, 3)$ 、 $B(4, 1)$ 、 $C(2, 5)$ 。写出 $\triangle ABC$ 关于 $x$ 轴对称后的顶点坐标。

\textbf{解：}
$A' = (1, -3)$ ， $B' = (4, -1)$ ， $C' = (2, -5)$ 。

\textbf{例 3.} 点 $P(a, b)$ 关于 $y$ 轴的对称点为 $Q$ ， $Q$ 关于 $x$ 轴的对称点为 $R$ 。求 $R$ 的坐标。 $R$ 与 $P$ 关于什么对称？

\textbf{解：}
$P(a, b) \xrightarrow{y\text{轴}} Q(-a, b) \xrightarrow{x\text{轴}} R(-a, -b)$ 。

$R(-a, -b)$ 就是 $P(a, b)$ 关于原点的对称点。

所以"先关于 $y$ 轴对称再关于 $x$ 轴对称"等价于"关于原点对称"。

\end{workedexamples}

\begin{keytakeaway}


\begin{itemize}
  \item $x$ 轴对称变 $y$ ， $y$ 轴对称变 $x$ ，原点对称都变。
  \item 两次轴对称（ $x$ 轴 $+$ $y$ 轴） $=$ 一次中心对称（关于原点）。
\end{itemize}


\end{keytakeaway}

\begin{practice}


\begin{enumerate}
  \item 点 $M(-5, 2)$ 关于 $x$ 轴的对称点是_____，关于 $y$ 轴的对称点是_____，关于原点的对称点是_____。
  \item 直线 $y = x + 2$ 关于 $x$ 轴对称后的方程是什么？
\end{enumerate}


\end{practice}



\section*{两点之间的距离公式}


\begin{explore}


两座城市在地图上的坐标分别是 $(2, 3)$ 和 $(5, 7)$ 。它们之间的"直线距离"是多少格？直觉告诉我们可以用勾股定理来算——水平差 $3$ 格、竖直差 $4$ 格，斜线距离就是 $5$ 格。

\end{explore}

\begin{understand}


\textbf{两点之间的距离公式}：

设 $A(x_1, y_1)$ 和 $B(x_2, y_2)$ ，则

\[
AB = \sqrt{(x_2 - x_1)^2 + (y_2 - y_1)^2}
\]


\begin{figure}[htbp]
\centering
\begin{tikzpicture}[every node/.style={font=\small}, scale=0.9]
  \coordinate (A) at (0,0);
  \coordinate (B) at (4,3);
  \coordinate (C) at (4,0);
  % Draw the right triangle
  \draw[main line] (A)--(B);
  \draw[aux line] (A)--(C)--(B);
  % Right angle at C
  \draw (3.75,0)--(3.75,0.25)--(4,0.25);
  % Points
  \fill[draw=red!70!black, fill=red!70!black] (A) circle (2.5pt)
    node[below left] {$A(x_1,y_1)$};
  \fill[draw=blue!70!black, fill=blue!70!black] (B) circle (2.5pt)
    node[above right] {$B(x_2,y_2)$};
  \fill (C) circle (2pt) node[below right] {};
  % Labels
  \node[below] at (2,0) {$|x_2-x_1|$};
  \node[right] at (4,1.5) {$|y_2-y_1|$};
  \node[above left, rotate={atan2(3,4)}] at (2,1.5)
    {$d = \sqrt{(x_2-x_1)^2+(y_2-y_1)^2}$};
\end{tikzpicture}
\caption{两点间距离公式：$|AB| = \sqrt{(x_2-x_1)^2+(y_2-y_1)^2}$}
\label{fig:distance-formula}
\end{figure}


推导：过 $A$ 、 $B$ 分别作坐标轴的平行线，构成直角三角形。水平距离 $= \lvert x_2 - x_1 \rvert$ ，竖直距离 $= \lvert y_2 - y_1 \rvert$ 。由勾股定理：

\[
AB = \sqrt{(x_2 - x_1)^2 + (y_2 - y_1)^2}
\]


特例：原点 $O(0, 0)$ 到点 $P(x, y)$ 的距离： $OP = \sqrt{x^2 + y^2}$ 。

\end{understand}

\begin{workedexamples}


\textbf{例 1.} 求 $A(1, 2)$ 与 $B(4, 6)$ 之间的距离。

\textbf{解：}
$AB = \sqrt{(4-1)^2 + (6-2)^2} = \sqrt{9 + 16} = \sqrt{25} = 5$ 。

\textbf{例 2.} 已知 $A(-1, 3)$ 、 $B(2, -1)$ ，求 $AB$ 的长度。

\textbf{解：}
$AB = \sqrt{(2-(-1))^2 + (-1-3)^2} = \sqrt{9 + 16} = \sqrt{25} = 5$ 。

\textbf{例 3.} 已知 $A(0, 0)$ 、 $B(6, 0)$ 、 $C(3, 4)$ 。证明 $\triangle ABC$ 是等腰三角形。

\textbf{解：}
$AB = \sqrt{(6-0)^2 + 0^2} = 6$ 。

$AC = \sqrt{3^2 + 4^2} = \sqrt{25} = 5$ 。

$BC = \sqrt{(6-3)^2 + (0-4)^2} = \sqrt{9 + 16} = 5$ 。

因为 $AC = BC = 5$ ，所以 $\triangle ABC$ 是等腰三角形。

\textbf{例 4.} 在 $x$ 轴上找一点 $P$ ，使 $PA = PB$ ，其中 $A(1, 3)$ 、 $B(5, 1)$ 。

\textbf{解：}
设 $P(x, 0)$ （在 $x$ 轴上）。

$PA = \sqrt{(x-1)^2 + 9}$ ， $PB = \sqrt{(x-5)^2 + 1}$ 。

由 $PA = PB$ ： $(x-1)^2 + 9 = (x-5)^2 + 1$ 。

$x^2 - 2x + 1 + 9 = x^2 - 10x + 25 + 1$

$-2x + 10 = -10x + 26$

$8x = 16$

$x = 2$

所以 $P(2, 0)$ 。

\end{workedexamples}

\begin{keytakeaway}


\begin{itemize}
  \item 距离公式 $= \sqrt{(\Delta x)^2 + (\Delta y)^2}$ ，本质是勾股定理。
  \item 距离公式可以用于判断三角形类型、求线段长度等。
\end{itemize}


\end{keytakeaway}

\begin{practice}


\begin{enumerate}
  \item 求 $A(-2, 1)$ 与 $B(4, -3)$ 之间的距离。
  \item 已知 $A(0, 0)$ 、 $B(4, 3)$ 、 $C(0, 5)$ ，判断 $\triangle ABC$ 的形状（按边分类）。
  \item 在 $y$ 轴上找一点 $P$ ，使 $P$ 到 $A(3, 4)$ 的距离为 $5$ 。
\end{enumerate}


\end{practice}



\section*{中点坐标公式}


\begin{explore}


两个好朋友住在城市的不同地方，他们想找一个"正中间"的地方见面。在坐标地图上，这个"正中间"的位置就是两点的中点。

\end{explore}

\begin{understand}


\textbf{中点坐标公式}：

设 $A(x_1, y_1)$ 和 $B(x_2, y_2)$ 的中点为 $M$ ，则

\[
M = \left(\frac{x_1 + x_2}{2}, \frac{y_1 + y_2}{2}\right)
\]


理解：中点的横坐标是两端点横坐标的平均值，纵坐标是两端点纵坐标的平均值。

\end{understand}

\begin{workedexamples}


\textbf{例 1.} 求 $A(2, 6)$ 与 $B(8, 4)$ 的中点坐标。

\textbf{解：}
$M = \left(\dfrac{2+8}{2}, \dfrac{6+4}{2}\right) = (5, 5)$ 。

\textbf{例 2.} 线段 $AB$ 的中点为 $M(3, -1)$ ， $A(1, 2)$ 。求 $B$ 的坐标。

\textbf{解：}
$\dfrac{1 + x_B}{2} = 3$ ， $x_B = 5$ 。

$\dfrac{2 + y_B}{2} = -1$ ， $y_B = -4$ 。

$B = (5, -4)$ 。

\textbf{例 3.} 平行四边形 $ABCD$ 的三个顶点为 $A(0, 0)$ 、 $B(4, 0)$ 、 $C(5, 3)$ 。求 $D$ 的坐标。

\textbf{解：}
平行四边形的对角线互相平分，所以 $AC$ 和 $BD$ 的中点相同。

$AC$ 的中点 $= \left(\dfrac{0+5}{2}, \dfrac{0+3}{2}\right) = \left(\dfrac{5}{2}, \dfrac{3}{2}\right)$ 。

$BD$ 的中点也是 $\left(\dfrac{5}{2}, \dfrac{3}{2}\right)$ ：

$\dfrac{4 + x_D}{2} = \dfrac{5}{2}$ ， $x_D = 1$ 。

$\dfrac{0 + y_D}{2} = \dfrac{3}{2}$ ， $y_D = 3$ 。

$D = (1, 3)$ 。

\textbf{例 4.} 已知 $\triangle ABC$ 的三个顶点为 $A(0, 6)$ 、 $B(-4, 0)$ 、 $C(4, 0)$ 。求 $BC$ 边上的中线 $AM$ 的长度。

\textbf{解：}
$M$ 是 $BC$ 的中点： $M = \left(\dfrac{-4+4}{2}, \dfrac{0+0}{2}\right) = (0, 0)$ 。

$AM = \sqrt{(0-0)^2 + (0-6)^2} = 6$ 。

\end{workedexamples}

\begin{keytakeaway}


\begin{itemize}
  \item 中点坐标 $= \left(\dfrac{x_1+x_2}{2}, \dfrac{y_1+y_2}{2}\right)$ ，即两端点坐标的平均值。
  \item 中点公式常与距离公式配合使用。
  \item 利用对角线中点相同可求平行四边形的顶点。
\end{itemize}


\end{keytakeaway}

\begin{practice}


\begin{enumerate}
  \item 求 $A(-3, 5)$ 与 $B(7, -1)$ 的中点坐标。
  \item 线段 $AB$ 的中点为 $M(4, 3)$ ， $A(2, 7)$ ，求 $B$ 的坐标。
  \item $\triangle ABC$ 中， $A(2, 4)$ 、 $B(6, 0)$ 、 $C(-2, -2)$ 。求三条中线的交点（重心）坐标。（提示：重心坐标 $= \left(\dfrac{x_A+x_B+x_C}{3}, \dfrac{y_A+y_B+y_C}{3}\right)$ 。）
\end{enumerate}


\end{practice}



\begin{answer}


\begin{enumerate}
  \item $A(5, -3)$ ：第四象限。 $B(-4, 0)$ ： $x$ 轴上（负半轴）。 $C(-2, -7)$ ：第三象限。 $D(0, 0)$ ：原点。
\end{enumerate}


\begin{enumerate}
  \item 在第二象限： $x < 0$ 且 $y > 0$ 。 $2a - 4 < 0$ ， $a < 2$ ； $a + 1 > 0$ ， $a > -1$ 。所以 $-1 < a < 2$ 。
\end{enumerate}


\begin{enumerate}
  \item 略（在坐标纸上标点即可）。
\end{enumerate}


\begin{enumerate}
  \item $OA = 6$ （沿 $x$ 轴），等腰直角三角形中 $OB = OA = 6$ ， $B$ 在 $y$ 轴正半轴上，所以 $B = (0, 6)$ 。
\end{enumerate}


\begin{enumerate}
  \item $(-3 + 5, 4 - 3) = (2, 1)$ 。
\end{enumerate}


\begin{enumerate}
  \item $A' = (1+3, 2) = (4, 2)$ ， $B' = (5+3, 6) = (8, 6)$ 。原线段长 $= \sqrt{16+16} = 4\sqrt{2}$ ，新线段长也为 $4\sqrt{2}$ （平移不改变长度）。
\end{enumerate}


\begin{enumerate}
  \item 关于 $x$ 轴： $(-5, -2)$ 。关于 $y$ 轴： $(5, 2)$ 。关于原点： $(5, -2)$ 。
\end{enumerate}


\begin{enumerate}
  \item 关于 $x$ 轴对称： $y \to -y$ 。原式 $y = x + 2$ 变为 $-y = x + 2$ ，即 $y = -x - 2$ 。
\end{enumerate}


\begin{enumerate}
  \item $AB = \sqrt{(4-(-2))^2 + (-3-1)^2} = \sqrt{36 + 16} = \sqrt{52} = 2\sqrt{13}$ 。
\end{enumerate}


\begin{enumerate}
  \item $AB = \sqrt{16+9} = 5$ ， $AC = \sqrt{0+25} = 5$ ， $BC = \sqrt{16+4} = \sqrt{20} = 2\sqrt{5}$ 。 $AB = AC = 5$ ，等腰三角形。
\end{enumerate}


\begin{enumerate}
  \item 设 $P(0, y)$ 。 $PA = \sqrt{9 + (y-4)^2} = 5$ 。 $9 + (y-4)^2 = 25$ ， $(y-4)^2 = 16$ ， $y - 4 = \pm 4$ 。 $y = 8$ 或 $y = 0$ 。两个点： $P(0, 8)$ 或 $P(0, 0)$ 。
\end{enumerate}


\begin{enumerate}
  \item $M = \left(\dfrac{-3+7}{2}, \dfrac{5+(-1)}{2}\right) = (2, 2)$ 。
\end{enumerate}


\begin{enumerate}
  \item $\dfrac{2+x_B}{2} = 4$ ， $x_B = 6$ 。 $\dfrac{7+y_B}{2} = 3$ ， $y_B = -1$ 。 $B = (6, -1)$ 。
\end{enumerate}


\begin{enumerate}
  \item 重心 $= \left(\dfrac{2+6+(-2)}{3}, \dfrac{4+0+(-2)}{3}\right) = \left(2, \dfrac{2}{3}\right)$ 。
\end{enumerate}


\end{answer}
